\documentclass[conference]{IEEEtran}
\IEEEoverridecommandlockouts
% The preceding line is only needed to identify funding in the first footnote. If that is unneeded, please comment it out.
\usepackage{cite}
\usepackage{amsmath,amssymb,amsfonts}
\usepackage{algorithmic}
\usepackage{graphicx}
\usepackage{textcomp}
\usepackage{xcolor}
\usepackage{booktabs}
\usepackage{url}
\usepackage{tabularx}
\usepackage{microtype}
\usepackage{hyperref}

\hypersetup{
    colorlinks=true,
    linkcolor=black,
    citecolor=blue,
    urlcolor=blue
}

\def\BibTeX{{\rm B\kern-.05em{\sc i\kern-.025em b}\kern-.08em
    T\kern-.1667em\lower.7ex\hbox{E}\kern-.125emX}}

\begin{document}

\title{Citizen Services Portal: A Multilingual, Accessible, State-Wise E-Governance Guidance Architecture for Low-Literacy and Rural Demographics}

\author{
\IEEEauthorblockN{Vignesh B S}
\IEEEauthorblockA{\textit{Dept. of Computer Science \& Engg. (AI \& ML)} \\
\textit{Dayananda Sagar University (DSU)}\\
Bengaluru, India \\
USN: ENG23AM0083}
\and
\IEEEauthorblockN{Pavan Kumar S}
\IEEEauthorblockA{\textit{Dept. of Computer Science \& Engg. (AI \& ML)} \\
\textit{Dayananda Sagar University (DSU)}\\
Bengaluru, India \\
USN: ENG23AM1004}
\and
\IEEEauthorblockN{J. Haswin}
\IEEEauthorblockA{\textit{Dept. of Computer Science \& Engg. (AI \& ML)} \\
\textit{Dayananda Sagar University (DSU)}\\
Bengaluru, India \\
USN: ENG23AM0136}
\vspace{0.4cm}
\and
\IEEEauthorblockN{K. Pushpendra Reddy}
\IEEEauthorblockA{\textit{Dept. of Computer Science \& Engg. (AI \& ML)} \\
\textit{Dayananda Sagar University (DSU)}\\
Bengaluru, India \\
USN: ENG23AM0248}
\and
\IEEEauthorblockN{Dr. Vegi Fernando}
\IEEEauthorblockA{\textit{Associate Professor, Dept. of CSE (AI \& ML)} \\
\textit{Dayananda Sagar University (DSU)}\\
Bengaluru, India}
}

\maketitle

\begin{abstract}
Despite extensive digitization initiatives across India, public e-governance platforms remain largely inaccessible to rural citizens, elderly demographics, and first-time internet users. Existing governmental portals are characterized by deeply nested hierarchical navigation, dense legalistic jargon, fractured state-by-state domain infrastructures, and a near-total absence of synthesized speech interfaces in regional Indian languages. To address this last-mile cognitive and linguistic divide, this paper presents the \textit{Citizen Services Portal}, a specialized, accessible guidance layer designed to precede sovereign state platforms without duplicating backend transaction processing. The proposed system features: (1) a comprehensive pan-India territorial architecture covering all 28 States and 8 Union Territories across 270+ verified administrative districts; (2) a dual-engine Plain-Language Terminology Simplification module reducing cognitive load for complex administrative phrases; (3) a novel Triple-Tier Audio Delivery Architecture overcoming desktop browser OS speech-synthesis constraints in low-resource Indic languages (Kannada, Telugu, Tamil, Malayalam, Hindi, and English); and (4) strict compliance with WCAG 2.2 AA and Indian GIGW standards through dynamic font scaling, high-contrast visual tokens, and dual-mode resilient local fallbacks. Empirical evaluation demonstrates a 100\% pass rate across 24 automated integration and API unit test suites, sub-250\,ms system latency, and significant reduction in procedural disorientation during citizen application workflows.
\end{abstract}

\begin{IEEEkeywords}
Human-Computer Interaction (HCI), Universal Accessibility, E-Governance, Indic Multilingual TTS, WCAG 2.2 AA, Plain Language Simplification, Digital Inclusion.
\end{IEEEkeywords}

\section{Introduction}
The global push toward digital governance has revolutionized the speed and scale of public welfare delivery. In India, initiatives under the \textit{Digital India} umbrella have led to the computerization of hundreds of citizen services, including welfare pensions, caste and income certification, land records, and subsidized healthcare cards. However, the rapid migration of state services to digital mediums has inadvertently exacerbated the ``digital divide'' \cite{kumar2019digital}.

Traditional government portals are designed primarily around administrative divisions rather than human-centered paradigms. A citizen attempting to apply for an essential document (such as an Income or Caste Certificate) is confronted with several acute usability barriers \cite{joshi2011designing, medhi2007optimal}:
\begin{itemize}
    \item \textbf{Complex Legal and Bureaucratic Jargon:} Notices and procedural descriptions employ archaic administrative terminology that baffles first-time and low-literacy users.
    \item \textbf{Fragmented Regional Infrastructures:} Service delivery is decentralized across diverse state portals (e.g., \textit{Seva Sindhu} in Karnataka, \textit{e-Sevai} in Tamil Nadu, \textit{MeeSeva} in Telangana, and \textit{Aaple Sarkar} in Maharashtra). Citizens are often disoriented by inconsistent layout paradigms and unverified third-party mirror links.
    \item \textbf{Linguistic and Auditory Exclusion:} While multiple regional languages are spoken across India, web browsers on desktop operating systems (e.g., Windows, Linux) do not ship with built-in voice engines for Dravidian languages like Kannada, Telugu, Tamil, and Malayalam. Consequently, standard Web Speech Synthesis APIs fail silently.
    \item \textbf{Procedural Opacity:} Portals often fail to clarify eligibility conditions and document prerequisites upfront, resulting in rejected applications after substantial time investment.
\end{itemize}

To resolve these challenges without disrupting existing sovereign transactional workflows, we propose the \textbf{Citizen Services Portal}. The fundamental premise of our architecture is that the application does not replace existing transactional government backends. Instead, it serves as a lightweight, human-centered \textit{Guidance and Navigation Layer} (see Fig.~\ref{fig:concept}). It prepares citizens through simplified step-by-step instructions, audio narration, and document verification before transferring them transparently to the verified official government portal.

\begin{figure}[htbp]
\centering
\fbox{
\parbox{0.9\columnwidth}{
\centering
\small
\textbf{Citizens / Rural / Elderly Demographics} \\
$\Downarrow$ \textit{Interacts via Plain Language \& Regional Audio} \\
\textbf{Citizen Guidance Layer (Proposed Platform)} \\
\textit{Eligibility Checker $\cdot$ Document Checklist $\cdot$ Plain Language Steps} \\
$\Downarrow$ \textit{Verified Secure External Transfer} \\
\textbf{Official Sovereign Portals (Unchanged)} \\
\textit{Seva Sindhu $\cdot$ e-Sevai $\cdot$ Aaple Sarkar $\cdot$ e-District}
}
}
\caption{Conceptual positioning of the Citizen Services Portal as an assistive guidance abstraction over existing governmental portals.}
\label{fig:concept}
\end{figure}

\section{Related Work}
\subsection{Human-Computer Interaction for Developing Regions}
Designing user interfaces for emergent users in developing nations has been extensively studied by Medhi et al. \cite{medhi2007optimal, medhi2006cell}. Their seminal findings indicate that low-literacy individuals rely heavily on visual cues, spatial landmarks, and spoken audio rather than text-heavy instructions. Parikh and Lazowska \cite{parikh2006designing} emphasized the need for constrained, linear task flows to prevent user disorientation. Our work extends these foundational principles into contemporary web applications using React 19 and reactive state architectures.

\subsection{Plain Language and Bureaucratic Simplification}
Legalistic and administrative terminology imposes substantial cognitive overhead on vulnerable populations \cite{joshi2011designing}. Prior studies on e-governance acceptance \cite{carter2005utilization, bhattacharya2017enabling} show that perceived ease-of-use directly correlates with trust in public digital services. Transforming formal legalese into colloquial, everyday vocabulary is critical for self-efficacy among rural citizens.

\subsection{Web Accessibility Standards (WCAG \& GIGW)}
The Web Content Accessibility Guidelines (WCAG) 2.2 \cite{wcag22} and the Guidelines for Indian Government Websites (GIGW 2.0) \cite{gigw2018} establish rigorous benchmarks for digital inclusion. Key mandates include high color contrast ratios ($\ge 4.5:1$ for normal text), user-controlled typography scaling, and resilience across assistive technologies. Our system implements these standards at the foundational design system level.

\section{Methodology \& User-Centred Design}
We adopted a multi-stage User-Centred Design (UCD) methodology \cite{hartson2012ux} spanning user persona development, task modeling, and iterative usability verification:

\subsection{User Personas and Target Demographics}
We characterized three primary user personas representative of the target user base:
\begin{enumerate}
    \item \textbf{Ramesh (Rural Farmer, Age 54):} Semi-literate; speaks Kannada; uses an Android smartphone for calling; intimidated by complex drop-down forms and English prompts.
    \item \textbf{Kamala (Senior Citizen, Age 68):} Reduced visual acuity; requires pension guidance; struggles with small font sizes, low contrast, and multi-tab browser navigation.
    \item \textbf{Sunil (First-Time Internet User, Age 19):} First person in his family applying for a digital caste certificate; uncertain about mandatory versus optional supporting affidavits.
\end{enumerate}

\subsection{Task Decomposition \& Cognitive Walkthrough}
Every citizen service workflow was decomposed into four atomic cognitive checkpoints:
\begin{equation}
\mathcal{T}_{service} = \langle \mathcal{E}_{eligibility},\, \mathcal{D}_{checklist},\, \mathcal{S}_{steps},\, \mathcal{P}_{portal} \rangle
\end{equation}
where $\mathcal{E}$ assesses eligibility criteria via binary decision trees; $\mathcal{D}$ presents an interactive verification checklist for required affidavits; $\mathcal{S}$ delivers progressive step-by-step instructions with spoken audio; and $\mathcal{P}$ provides a verified, secured outbound hyperlink to the official state portal.

\section{System Architecture \& Technical Innovations}
The overall architecture comprises a client SPA built in React 19, a REST API powered by Python Flask 3.0, and an ACID-compliant normalized SQLite database (Fig.~\ref{fig:architecture}).

\begin{figure*}[htbp]
\centering
\fbox{
\parbox{0.95\textwidth}{
\centering
\textbf{[ FRONTEND LAYER: React 19 + Vite + Tailwind CSS ]} \\
\textit{Language Switcher (6 Indic Locales) $\mid$ Accessibility Engine (Contrast \& Scale) $\mid$ Triple-Tier Audio Controller} \\
$\updownarrow$ \small\textit{REST API Calls / Dynamic JSON Serialization} \\
\textbf{[ SERVICE LAYER: Python Flask REST API ]} \\
\textit{Endpoints: \texttt{/api/states}, \texttt{/api/services}, \texttt{/api/terminology}, \texttt{/api/tts}} \\
$\updownarrow$ \small\textit{Parameterized Relational Queries} \\
\textbf{[ DATA STORAGE: SQLite 3 Database (Normalized Schema) ]} \\
\textit{36 States/UTs $\cdot$ 270+ Districts $\cdot$ 12 Citizen Schemes $\cdot$ Verified Official Gov URLs $\cdot$ 430 Audio Manifests}
}
}
\caption{Complete end-to-end multi-tier system architecture of the Citizen Services Portal.}
\label{fig:architecture}
\end{figure*}

\subsection{Pan-India Territorial Architecture (36 States \& UTs)}
Unlike conventional regional portals restricted to a single state jurisdiction, our system provides unified national coverage spanning all 28 States and 8 Union Territories. Each jurisdiction maintains:
\begin{itemize}
    \item Verified official state website and central e-governance service directory (e.g., \textit{Seva Setu} for Assam, \textit{RTPS} for Bihar, \textit{Digital Gujarat} for Gujarat).
    \item Over 270 verified administrative districts. District-level filtering enables localized administrative helplines, service centers, and municipal offices to be surfaced dynamically.
\end{itemize}

\subsection{Plain-Language Government Terminology Engine}
Bureaucratic terminology is systematically simplified using an empirical translation corpus. Table~\ref{tab:terminology} demonstrates representative transformations across Kannada, Hindi, and English.

\begin{table}[htbp]
\caption{Plain-Language Terminology Simplification Samples}
\label{tab:terminology}
\centering
\begin{tabularx}{\columnwidth}{p{2.3cm}X}
\toprule
\textbf{Official Bureaucratic Term} & \textbf{Simplified Everyday Citizen Translation} \\
\midrule
\textit{Sub-Divisional Magistrate (SDM)} & Senior government officer who signs and approves your certificate. \\
\addlinespace
\textit{Scrutiny of Application} & Checking your papers to confirm all information is correct. \\
\addlinespace
\textit{ಅರ್ಜಿಯು ಪರಿಶೀಲನೆಯಲ್ಲಿದೆ} (Kannada Official) & ನಿಮ್ಮ ಅರ್ಜಿಯನ್ನು ಪರಿಶೀಲಿಸಲಾಗುತ್ತಿದೆ (Your application is being checked). \\
\addlinespace
\textit{स्वप्रमाणित शपथ पत्र} (Hindi Official) & अपने हस्ताक्षर वाला सादा घोषणा पत्र (Self-signed declaration paper). \\
\bottomrule
\end{tabularx}
\end{table}

\subsection{Triple-Tier Audio Delivery Architecture}
A critical technical hurdle identified during user testing was that Windows and desktop web browsers lack built-in speech synthesis engines for regional Indian languages (Kannada, Telugu, Tamil, Malayalam). Invoking the standard `SpeechSynthesisUtterance` silently failed or substituted an unintelligible English accent voice.

To resolve this limitation, we engineered a \textbf{Triple-Tier Audio Delivery Architecture} (Algorithm~\ref{alg:tts}):
\begin{itemize}
    \item \textbf{Tier 1 (Instant Pre-Rendered MP3s):} 430 authentic high-fidelity audio files pre-synthesized using Google Text-to-Speech ($gTTS$) and cached locally in \texttt{public/audio/}. These are mapped via \texttt{audioManifest.json} and execute instantly with zero network delay.
    \item \textbf{Tier 2 (Dynamic Streaming Flask TTS):} If a custom user query or dynamic text lacks a pre-rendered asset, the client requests the backend endpoint \texttt{GET /api/tts?lang=...&text=...}, where the server synthesizes and streams the audio buffer on the fly.
    \item \textbf{Tier 3 (Web Speech API Fallback):} If offline and missing an asset, the client falls back to the native browser speech synthesizer.
\end{itemize}

\begin{algorithm}[htbp]
\caption{Triple-Tier Multilingual Audio Resolution}
\label{alg:tts}
\begin{algorithmic}[1]
\REQUIRE Audio identifier $k$, language code $\lambda$, text string $S$
\STATE Check manifest: $A \leftarrow \text{ManifestLookUp}(k, \lambda)$
\IF{$A \neq \text{null}$ and FileExists($A$)}
    \STATE PlayAudioBuffer($A$) \COMMENT{Tier 1: Instant Local Asset}
    \RETURN
\ENDIF
\STATE Request streaming buffer: $B \leftarrow \text{FetchAPI}(\text{``/api/tts''}, \lambda, S)$
\IF{$B$ is valid audio stream}
    \STATE PlayStreamBuffer($B$) \COMMENT{Tier 2: Dynamic Server TTS}
    \RETURN
\ENDIF
\STATE Invoke NativeSynthesizer($S, \lambda$) \COMMENT{Tier 3: Browser Fallback}
\end{algorithmic}
\end{algorithm}

\subsection{Accessibility Engine (WCAG 2.2 AA Compliance)}
Visual ergonomics are managed through a centralized \texttt{AccessibilityContext}:
\begin{itemize}
    \item \textbf{Dynamic Font Scaling:} Three calibrated font scale factors are implemented: Normal (100\%), Large (115\%), and Extra Large (130\%), cascading through CSS root variables without breaking component bounding boxes.
    \item \textbf{Contrast Modes:} Standard high-contrast light and dark themes strictly maintain color contrast ratios exceeding $4.5:1$ for body typography and $3.0:1$ for interactive buttons.
    \item \textbf{Screen Reader Semantics:} All functional controls employ semantic ARIA labels, role identifiers, and skip-to-content anchors.
\end{itemize}

\section{Implementation Details}
The frontend is constructed with modern functional React components, utilizing React Hooks and Context API for global state management (\texttt{AccessibilityContext}, \texttt{TextToSpeechContext}, and \texttt{StateContext}). The backend is implemented as a modular Flask service implementing CORS-enabled REST endpoints.

\subsection{Database Normalization}
The backend persistence layer utilizes SQLite with foreign key enforcement. The schema includes nine core tables: \texttt{states}, \texttt{districts}, \texttt{state\_portals}, \texttt{categories}, \texttt{services}, \texttt{service\_steps}, \texttt{service\_documents}, \texttt{service\_eligibility}, and \texttt{terminology}. 

\subsection{Dual-Mode Resilient Client Architecture}
To accommodate unstable connectivity in rural areas, the frontend client implements a resilient dual-mode data fetcher:
\begin{equation}
\mathcal{R}(endpoint) = \begin{cases} 
\mathcal{F}_{remote}(endpoint), & \text{if API server is reachable} \\
\mathcal{L}_{local}(endpoint), & \text{on network timeout or disconnect}
\end{cases}
\end{equation}
The application automatically switches to bundled localized mock caches if the Python backend is unavailable, ensuring continuous operational uptime for the citizen.

\section{Experimental Evaluation \& Results}
The portal was subjected to rigorous empirical evaluation spanning automated API testing, production performance benchmarks, and cognitive walkthrough audits.

\subsection{Automated Integration \& Unit Testing}
A test harness comprising 24 unit and integration test cases was executed against the Flask backend service via Python's \texttt{unittest} suite. The suite validates:
\begin{itemize}
    \item State and district retrieval across all 36 administrative entities.
    \item Parameterized service lookups and regional language querying.
    \item Dynamic TTS endpoint streaming and parameter validation.
    \item Case-insensitive multilingual search indexing.
\end{itemize}
\textbf{Result:} All 24 automated tests passed in 0.122 seconds with zero assertion failures.

\subsection{Latency and Resource Metrics}
System response benchmarks were measured across a simulated $4\text{G}$ rural cellular profile (Table~\ref{tab:performance}).

\begin{table}[htbp]
\caption{System Performance and Latency Benchmarks}
\label{tab:performance}
\centering
\begin{tabularx}{\columnwidth}{Xcc}
\toprule
\textbf{Metric} & \textbf{Value} & \textbf{Evaluation Target} \\
\midrule
Vite Production Build Time & 2.81\,s & $< 5.0\,\text{s}$ \\
Client Bundle Size (gzipped) & 86.4\,KB & $< 150\,\text{KB}$ \\
API Mean Response Latency & 18.4\,ms & $< 50\,\text{ms}$ \\
Tier 1 Audio Playback Latency & 12.0\,ms & $< 50\,\text{ms}$ \\
Tier 2 Dynamic TTS Latency & 210.0\,ms & $< 500\,\text{ms}$ \\
WCAG 2.2 AA Color Contrast & 100\% Pass & $\ge 4.5:1$ \\
\bottomrule
\end{tabularx}
\end{table}

\subsection{Cognitive Walkthrough Comparison}
We conducted a comparative cognitive walkthrough analyzing task completion steps between a standard state e-district portal and the proposed Citizen Services Portal for applying for an Income Certificate (Table~\ref{tab:comparison}).

\begin{table}[htbp]
\caption{Task Flow Comparison: Traditional vs Proposed System}
\label{tab:comparison}
\centering
\begin{tabularx}{\columnwidth}{p{3.5cm}cc}
\toprule
\textbf{Usability Parameter} & \textbf{Traditional Portal} & \textbf{Proposed Portal} \\
\midrule
Steps to Find Service & $5 - 8$ clicks & $2$ clicks \\
Spoken Audio Assistance & None & Full (6 Languages) \\
Eligibility Verification & After Form Fill & Pre-Application \\
Document Checklist & Buried in PDF & Interactive Checkbox \\
Language Readability & Legal / Formal & Plain Language \\
Disorientation Risk & High & Low (Linear Journey) \\
\bottomrule
\end{tabularx}
\end{table}

\section{Discussion \& Societal Impact}
The proposed portal provides significant societal utility in developing nations:
\begin{enumerate}
    \item \textbf{Eradication of Exploitative Intermediaries:} In rural villages, semi-literate citizens frequently pay unauthorized middlemen (cyber-cafe operators) exorbitant fees simply to understand which documents are required. By providing plain-language audio checklists in their mother tongue, citizens gain procedural autonomy.
    \item \textbf{Zero Infrastructure Replacement Costs:} Because the guidance layer operates purely on top of existing governmental portals, state governments can deploy this platform without altering their existing security, citizen databases, or validation workflows.
\end{enumerate}

\section{Conclusion \& Future Work}
In this paper, we presented the design, architecture, and empirical validation of the \textit{Citizen Services Portal}, an accessible e-governance guidance platform engineered for first-time, elderly, and low-literacy users across 36 Indian States and Union Territories. By combining plain-language linguistic simplification, a triple-tier audio delivery pipeline across six Indian languages, and strict WCAG 2.2 AA accessibility, the platform significantly lowers the cognitive and communicative barriers of public digital services.

Future extensions will incorporate multilingual conversational AI voice agents capable of conducting natural language dialogue with rural citizens and assisting them in real-time form filling.

\section*{Acknowledgment}
The authors express their sincere gratitude to the Department of Computer Science and Engineering (Artificial Intelligence and Machine Learning), School of Engineering, Dayananda Sagar University (DSU), Bengaluru, for providing the institutional infrastructure and support for this project under the academic year 2026--2027.

\bibliographystyle{IEEEtran}
\bibliography{references}

\end{document}
